% Part: methods
% Chapter: proofs
% Section: proof-by-contradiction

\documentclass[../../../include/open-logic-section]{subfiles}

\begin{document}

\olfileid{mth}{prf}{con}

\olsection{વિરોધાભાસથી સાબિતી}

સૌપ્રથમ, વિરોધાભાસથી સાબિતી નકારાત્મક દાવા સાબિત કરવા વપરાતો
અનુમાનનો નમૂનો છે. ધારો કે કોઈ દાવો~$p$ \emph{ખોટો} છે,
એટલે~$\lnot p$ છે, તે બતાવવું છે. સૌથી આશાસ્પદ રીત આ છે:
(a) $p$ સાચું છે એમ ધારો અને (b) બતાવો કે આ પૂર્વધારણાથી
ખોટી હોવાનું જાણીતું કંઈક નીકળે છે. ``ખોટી હોવાનું જાણીતું''
પરિણામ પોતે~$p$ સાથે અથવા તપાસાતા સમગ્ર દાવાની બીજી કોઈ
પૂર્વધારણા સાથે ટકરાઈ શકે. દાખલા તરીકે, ``જો $q$ તો $\lnot p$''
સાબિત કરવા $q$ સાચું ધારીને તેમાંથી~$\lnot
p$ સાબિત કરવું પડે. જો $\lnot p$ વિરોધાભાસથી સાબિત કરો,
તો~$q$ ઉપરાંત $p$ પણ ધારવું પડે. $p$ પરથી $\lnot q$
સાબિત થાય, તો~$p$ની પૂર્વધારણાથી તમારી બીજી પૂર્વધારણા~$q$
સાથે વિરોધાભાસ મળે છે, કારણ કે $q$ અને $\lnot q$ બંને
સાચાં હોઈ શકતાં નથી. અલબત્ત, વિરોધાભાસ મેળવવા સાબિતીમાં
બીજા અનુમાન-નમૂનાઓ અને વ્યાખ્યાઓ ઉકેલવાની રીત પણ વાપરવી
પડે. એક ઉદાહરણ જોઈએ.

\begin{prop}
  જો $A \subseteq B$ અને $B = \emptyset$, તો $A$માં કોઈ
  !!{element}s નથી.
\end{prop}

\begin{proof}
  $A \subseteq B$ અને $B = \emptyset$ ધારો. આપણે
  $A$માં કોઈ !!{element}s નથી તે બતાવવું છે.
  \begin{quote}
    આ શરતી દાવો હોવાથી પૂર્વપક્ષ ધારીને પરિણામ સાબિત
    કરવાનું છે. પરિણામ છે: $A$માં કોઈ !!{element}s નથી.
    વધુ સ્પષ્ટ કહીએ તો એવો~$x \in A$ છે એવું નથી.
  \end{quote}
  $A$માં કોઈ !!{element}s નથી ત્યારે અને ત્યારે જ જ્યારે
  એવો~$x$ છે કે $x \in A$ એવું નથી.
  \begin{quote}
    એટલે આપણે ખરેખર નકારાત્મક દાવો~$\lnot p$ સાબિત કરવો છે:
    એવો $x \in A$ છે એવું નથી. વિરોધાભાસથી સાબિતી કરવા
    અનુરૂપ હકારાત્મક દાવો~$p$, એટલે એવો $x \in A$ છે, એમ
    ધારીને તેમાંથી વિરોધાભાસ મેળવવો પડે. ``વિરોધાભાસ માટે
    ધારો કે'' અથવા ફક્ત ``ધારો કે એવું નથી'' લખીને આ રીત
    દર્શાવીએ છીએ અને પછી પૂર્વધારણા~$p$ કહીએ છીએ.
  \end{quote}
  ધારો કે એવું નથી: કોઈ $x \in A$ છે.
  \begin{quote}
    હવે આ નવી પૂર્વધારણાથી વિરોધાભાસ મેળવવો છે. બીજી બે
    પૂર્વધારણાઓ $A \subseteq B$ અને $B = \emptyset$ છે.
    પહેલી પરથી $x \in B$ મળે છે:
  \end{quote}
  $A \subseteq B$ હોવાથી $x \in B$.
  \begin{quote}
    પરંતુ $B = \emptyset$ હોવાથી $B$નું દરેક !!{element}
    (જેમ કે $x$)~$\emptyset$નું પણ !!a{element} હોવું જોઈએ.
  \end{quote}
  $B = \emptyset$ હોવાથી $x \in \emptyset$. આ વિરોધાભાસ છે,
  કારણ કે વ્યાખ્યા મુજબ $\emptyset$માં કોઈ !!{element}s નથી.
  \begin{quote}
    આથી સાબિતી પૂરી થાય છે: લીધેલી પૂર્વધારણાઓમાંથી જરૂરી
    વિરોધાભાસ મળ્યો છે, તેથી બધી પૂર્વધારણાઓ એકસાથે સાચી
    હોઈ શકતી નથી. પહેલી બે પૂર્વધારણાઓ ($A \subseteq B$ અને
    $B = \emptyset$) પર શંકા નથી, તેથી છેલ્લે લીધેલી
    પૂર્વધારણા (એવો $x \in A$ છે) ખોટી હોવી જોઈએ. વિગત
    જોઈએ તો તેને સ્પષ્ટ લખી શકીએ.
  \end{quote}
  તેથી એવો $x \in A$ છે એવી આપણી પૂર્વધારણા ખોટી છે;
  એટલે વિરોધાભાસથી સાબિતી મુજબ $A$માં કોઈ !!{element}s નથી.
\end{proof}

શાસ્ત્રીય તર્કમાં દરેક હકારાત્મક દાવો નકારાત્મક દાવાને સમકક્ષ
છે: $p$ ત્યારે અને ત્યારે જ જ્યારે $\lnot\lnot p$.
\sourcecorrection{OLMD-159}{મૂળમાં આ દ્વિનિષેધ સમકક્ષતા કોઈ
વ્યાપશરત વિના કહેવાઈ છે. તે શાસ્ત્રીય તર્કમાં માન્ય છે; સ્ફુરણાવાદી
તર્કમાં સામાન્ય રીતે નહીં.} તેથી વિરોધાભાસથી સાબિતી હકારાત્મક
દાવાને ``પરોક્ષ રીતે'' સ્થાપિત કરવા પણ વાપરી શકાય છે:~$p$
સાબિત કરવા તેને નકારાત્મક દાવો $\lnot\lnot p$ તરીકે વાંચો.
$\lnot p$ પરથી વિરોધાભાસ સાબિત થાય, તો વિરોધાભાસથી સાબિતી
મુજબ $\lnot\lnot p$ મળે અને તેથી~$p$ મળે.

છેલ્લા ઉદાહરણમાં આપણું લક્ષ્ય $A$માં કોઈ !!{element}s નથી એવો
નકારાત્મક દાવો હતો. તેથી વિરોધાભાસથી સાબિતી માટે લીધેલી
પૂર્વધારણા (એવો $x \in A$ છે) હકારાત્મક હતી. તેમાંથી વિચારવા
માટે કશુંક મળ્યું, એટલે કામચલાઉ $x \in A$; તેના વિશે આગળ
વિચારતાં $x \in \emptyset$ સુધી પહોંચ્યા.

હકારાત્મક દાવો પરોક્ષ રીતે સાબિત કરતાં વિરોધાભાસ માટેની
પૂર્વધારણા નકારાત્મક હોય છે. પરંતુ ઘણી વાર હકારાત્મક દાવાને
નકારાત્મક સ્વરૂપમાં અને નકારાત્મક દાવાને હકારાત્મક સ્વરૂપમાં
સહેલાઈથી કહી શકાય. જો છેલ્લા ઉદાહરણમાં ``$A$માં કોઈ
!!{element}s નથી'' એ નકારાત્મક પરિણામને બદલે ``$A = \emptyset$''
સાબિત કર્યું હોત તો સાબિતી લગભગ એવી જ હોત. ($=$ની વ્યાખ્યા
મુજબ ``$A = \emptyset$'' સાર્વત્રિક દાવો છે: તેને ઉકેલતાં
``~$A$નું દરેક !!{element},~$\emptyset$નું !!{element} છે
અને ઊલટું પણ'' મળે છે.) પરંતુ તે ``એવો $x \in A$ છે
એવું નથી'' એ નકારાત્મક દાવાને સમકક્ષ છે તે સહેલાઈથી જોઈ શકાય.

આથી ક્યારેક~$p$ને સીધું સાબિત કરવા કરતાં $\lnot p$ને
પૂર્વધારણા તરીકે વાપરવું સરળ હોય છે. સીધી સાબિતી એટલી જ સરળ
અથવા વધુ સરળ હોય (જેમ કે આગળનાં ઉદાહરણોમાં), તોપણ કેટલાક
લોકો પરોક્ષ રીતે આગળ વધવાનું પસંદ કરે છે. દ્વિનિષેધ ગૂંચવે તો
વિરોધાભાસથી કોઈ દાવાની સાબિતીને તેના \emph{વિપરીત} દાવામાંથી
વિરોધાભાસ મેળવવાની સાબિતી માનો. એટલે $\lnot
p$ની વિરોધાભાસથી સાબિતી એટલે~$p$ ધારીને વિરોધાભાસ મેળવવો;
અને~$p$ની વિરોધાભાસથી સાબિતી એટલે~$\lnot p$ ધારીને
વિરોધાભાસ મેળવવો.

\begin{prop}
$A \subseteq A \cup B$.
\end{prop}

\begin{proof}
  આપણે $A \subseteq A \cup B$ બતાવવું છે.
  \begin{quote}
    ઉપરથી આ હકારાત્મક દાવો છે: દરેક $x \in A$, $A \cup B$માં
    પણ છે. તેનો નિષેધ છે: કોઈ $x \in A$ એવો છે જે
    $\notin A \cup B$ છે. આ નકારેલા દાવાને ધારીને તેમાંથી
    વિરોધાભાસ મેળવી દાવો પરોક્ષ રીતે સાબિત કરી શકીએ.
    \end{quote}
  ધારો કે એવું નથી, એટલે $A \nsubseteq A \cup B$.
  \begin{quote}
    $A \subseteq A \cup B$ની વ્યાખ્યા આપણી પાસે છે: દરેક
    $x \in A$, $\in A \cup B$ પણ છે. $A \nsubseteq A \cup B$
    શું કહે છે તે સમજવા ઉકેલેલી વ્યાખ્યા પર થોડું મૂળભૂત તાર્કિક
    કામ કરવું પડે: દરેક $x \in A$~$\in
    A \cup B$ પણ છે એવું ખોટું હોય ત્યારે અને ત્યારે જ જ્યારે
    \emph{કોઈ}~$x \in A$ એવો છે જે $\notin A \cup B$ છે.
    \sourcecorrection{OLMD-160}{મૂળના આ વાક્યમાં અહીં અચાનક
    Cમાં સભ્ય નથી લખાયું છે; દાવાના નિષેધ મુજબ A યોગ Bમાં
    સભ્ય નથી હોવું જોઈએ.} (અહીં ખાસ સાવધ રહો: ``બધા $A$s
    એ $B$s છે'' જેવા દાવાનો નિષેધ ખોટી રીતે સમજવાથી
    વિદ્યાર્થીઓની ઘણી વિરોધાભાસી સાબિતીઓ ખોટી પડે છે.)
    બીજા શબ્દોમાં, $A \nsubseteq A \cup B$ ત્યારે અને ત્યારે જ
    જ્યારે એવો $x$ છે કે $x \in A$ અને $x \notin A \cup B$.
    પછીનું સરળ છે.
  \end{quote}
  તેથી એવો $x \in A$ છે કે $x \notin A \cup B$. $\cup$ની
  વ્યાખ્યા મુજબ $x \in A \cup B$ ત્યારે અને ત્યારે જ જ્યારે
  $x \in A$ અથવા $x \in
  B$. $x \in A$ હોવાથી $x \in A \cup B$ મળે છે. પરંતુ આ
  $x \notin A \cup B$ની પૂર્વધારણાથી વિરોધાભાસી છે.
\end{proof}

\begin{prob}
$A \cap B \subseteq A$ \emph{પરોક્ષ રીતે} સાબિત કરો.
\end{prob}

\begin{prop}
જો $A \subseteq B$ અને $B \subseteq C$, તો $A \subseteq C$.
\end{prop}

\begin{proof}
  $A \subseteq B$ અને $B \subseteq C$ ધારો. આપણે $A
  \subseteq C$ બતાવવું છે.
  \begin{quote}
    પરોક્ષ રીતે આગળ વધીએ: જે સાબિત કરવું છે તેનો નિષેધ ધારો.
  \end{quote}
  ધારો કે એવું નથી, એટલે $A \nsubseteq C$.
  \begin{quote}
    અગાઉની જેમ, $A \nsubseteq C$ ત્યારે અને ત્યારે જ જ્યારે
    દરેક $x \in A$, $\in C$ પણ છે એવું નથી; એટલે કોઈ
    $x \in A$ એવો છે જે $\notin C$ છે. ચિંતા ન કરો,
    અભ્યાસથી આવા નિષેધ ઉકેલવા વધારે વિચારવું નહીં પડે.
  \end{quote}
  બીજા શબ્દોમાં, એવો~$x$ છે કે $x \in A$ અને $x \notin C$.
  \begin{quote}
    હવે અહીંથી વિરોધાભાસ મેળવીએ. એ માટે બીજી બે પૂર્વધારણાઓ
    પણ વાપરવી પડશે.
  \end{quote}
  $A \subseteq B$ હોવાથી $x \in B$. $B \subseteq C$ હોવાથી
  $x \in
  C$. પરંતુ આ $x \notin C$ સાથે વિરોધાભાસ છે.
\end{proof}

\begin{prop}
જો $A \cup B = A \cap B$ તો $A = B$.
\end{prop}

\begin{proof}
  $A \cup B = A \cap B$ ધારો. આપણે $A = B$ બતાવવું છે.
  \begin{quote}
    હવે શરૂઆતની રીત જાણીતી છે:
  \end{quote}
  વિરોધાભાસ માટે $A \neq B$ ધારો.
  \begin{quote}
    વિરોધાભાસી સાબિતીની આપણી પૂર્વધારણા $A \neq
    B$ છે. $A = B$ ત્યારે અને ત્યારે જ જ્યારે $A \subseteq B$
    અને $B \subseteq A$; તેથી $A \neq B$ ત્યારે અને ત્યારે જ
    જ્યારે $A \nsubseteq B$ \emph{અથવા} $B \nsubseteq
    A$. (નિષેધનો વિચાર કરતી વખતે સાવધ રહેવું કેટલું અગત્યનું
    છે તે જુઓ!{}) આ વિયોજનમાંથી વિરોધાભાસ મેળવવા કિસ્સાવાર
    સાબિતી કરીએ: દરેક કિસ્સામાં વિરોધાભાસ મળે છે તે બતાવીએ.
  \end{quote}
  $A \neq B$ ત્યારે અને ત્યારે જ જ્યારે $A \nsubseteq B$ અથવા
  $B \nsubseteq A$. કિસ્સા અલગ પાડીએ.
  \begin{quote}
    પહેલા કિસ્સામાં $A \nsubseteq B$ ધારો; એટલે કોઈ $x$
    માટે $x \in A$ પણ $\notin B$. $A \cap B$ એ $A$ અને
    $B$ બંનેનાં સામાન્ય !!{element}sનો ગણ છે; તેથી કોઈ
    વસ્તુ એક ગણમાં ન હોય તો તે છેદગણમાં પણ નથી. $A \cup B$
    એ $A$ અને $B$નાં બધા સભ્યોનો ગણ છે, એટલે કોઈ એકમાં
    હોય તે યોગગણમાં પણ છે. તેથી $x \in A \cup B$ પણ
    $x \notin A \cap
    B$, અને આથી $A \cap B \neq A \cup B$.
  \end{quote}

  કિસ્સો 1: $A \nsubseteq B$. તો કોઈ $x$ માટે $x \in A$
  પણ $x \notin
  B$. $x \notin B$ હોવાથી $x \notin A \cap B$. $x \in A$
  હોવાથી $x \in A \cup B$. તેથી $A \cap B \neq A \cup B$,
  જે $A \cap B = A \cup B$ની પૂર્વધારણાથી વિરોધાભાસી છે.

  કિસ્સો 2: $B \nsubseteq A$. તો કોઈ $y$ માટે $y \in B$
  પણ $y \notin
  A$. અગાઉની જેમ $y \in A \cup B$ પરંતુ $y \notin A \cap B$;
  તેથી $A \cap B \neq A \cup B$, જે ફરી $A \cap B = A \cup
  B$થી વિરોધાભાસી છે.
\end{proof}

\end{document}
